% =========================================================
% 3.12 Register
% =========================================================

\section{রেজিস্টার
(\textenglish{Register})}

\subsection{রেজিস্টার কী?}

একাধিক \textenglish{Flip-Flop} একত্রে সংযুক্ত করে যে ডিজিটাল সার্কিটে
একাধিক বিটের বাইনারি তথ্য সাময়িকভাবে সংরক্ষণ করা যায়,
তাকে \textbf{\textenglish{Register}} বলে।

প্রতিটি \textenglish{Flip-Flop} সাধারণত একটি বিট তথ্য সংরক্ষণ করতে পারে।
তাই \(n\)-বিট তথ্য সংরক্ষণের জন্য সাধারণত \(n\)টি
\textenglish{Flip-Flop} প্রয়োজন হয়।

অর্থাৎ,

\[
\boxed{
n\text{টি }\text{\textenglish{Flip-Flop}}
\Rightarrow
n\text{-}\text{\textenglish{bit}}\text{ তথ্য সংরক্ষণ}
}
\]

উদাহরণস্বরূপ, একটি \(4\)-\textenglish{bit Register}-এ

\[
1011
\]

সংরক্ষণ করতে চারটি \textenglish{Flip-Flop} ব্যবহার করা যায়।


% =========================================================
% 3.12.2 Construction
% =========================================================

\subsection{রেজিস্টারের গঠন
(\textenglish{Construction of Register})}

রেজিস্টার সাধারণত একাধিক \textenglish{Flip-Flop} দিয়ে তৈরি করা হয়।
ডিজিটাল সিস্টেমে \textenglish{D Flip-Flop} ব্যবহার করে Register তৈরি করা
খুবই প্রচলিত।

একটি \(4\)-\textenglish{bit Register}-এ চারটি
\textenglish{Flip-Flop} ব্যবহার করলে প্রতিটি
\textenglish{Flip-Flop} একটি করে বিট সংরক্ষণ করে।

\[
\boxed{
FF_3+FF_2+FF_1+FF_0
=
4\text{-}\text{\textenglish{bit Register}}
}
\]

রেজিস্টারের প্রধান অংশগুলো হলো—

\begin{enumerate}[leftmargin=*,itemsep=2pt]

    \item \textbf{\textenglish{Flip-Flop}:}
    প্রতিটি বিট সংরক্ষণ করে।

    \item \textbf{\textenglish{Clock Input}:}
    কখন তথ্য গ্রহণ বা পরিবর্তন হবে তা নিয়ন্ত্রণ করে।

    \item \textbf{\textenglish{Data Input/Output}:}
    বাইনারি তথ্য গ্রহণ ও সরবরাহ করে।

    \item \textbf{\textenglish{Control Logic}:}
    প্রয়োজন অনুযায়ী তথ্য লোড, স্থানান্তর বা ধরে রাখার কাজ নিয়ন্ত্রণ করে।

\end{enumerate}


% =========================================================
% 3.12.3 Applications
% =========================================================

\subsection{রেজিস্টারের ব্যবহার}

রেজিস্টার বিভিন্ন ডিজিটাল ডিভাইসে ব্যবহৃত হয়। যেমন—

\begin{enumerate}[leftmargin=*,itemsep=2pt]

    \item \textenglish{CPU}-তে নির্দেশনা, ডেটা ও মধ্যবর্তী ফলাফল
    সাময়িকভাবে সংরক্ষণ করতে।

    \item বাইনারি তথ্য এক স্থান থেকে অন্য স্থানে স্থানান্তর করতে।

    \item \textenglish{Serial} ও \textenglish{Parallel Data}-এর মধ্যে
    রূপান্তর করতে।

    \item ডিজিটাল সার্কিটে তথ্য সাময়িকভাবে ধরে রাখতে।

    \item গাণিতিক ও লজিক্যাল অপারেশনের সময় তথ্য সংরক্ষণ করতে।

\end{enumerate}


% =========================================================
% 3.12.4 Shift Register
% =========================================================

\subsection{শিফট রেজিস্টার
(\textenglish{Shift Register})}

যে রেজিস্টারে \textenglish{Clock Pulse}-এর সাহায্যে সংরক্ষিত বাইনারি তথ্য
এক বিট করে বাম বা ডান দিকে স্থানান্তর করা যায়,
তাকে \textbf{\textenglish{Shift Register}} বলে।

\subsubsection*{গঠন ও কার্যপদ্ধতি}

\textenglish{Shift Register} সাধারণত একাধিক
\textenglish{D Flip-Flop} ধারাবাহিকভাবে সংযুক্ত করে তৈরি করা হয়।

একটি \textenglish{Flip-Flop}-এর \(Q\) আউটপুট পরবর্তী
\textenglish{Flip-Flop}-এর \(D\) ইনপুটের সঙ্গে যুক্ত থাকে।
সবগুলো \textenglish{Flip-Flop} একই \textenglish{Clock} সংকেত গ্রহণ করে।

প্রতিটি \textenglish{Clock Pulse}-এ সংরক্ষিত বিটগুলো এক ধাপ করে
পরবর্তী অবস্থানে স্থানান্তরিত হয়।


% =========================================================
% 4-bit Shift Register Diagram
% =========================================================

\subsubsection*{৪-বিট শিফট রেজিস্টারের গঠন
(\textenglish{4-bit Shift Register})}

\begin{center}
\begin{chThreeDiagram}
\begin{tikzpicture}[
    >=Stealth,
    line width=0.8pt,
    ff/.style={
        draw,
        rounded corners=2pt,
        fill=ICTLightBlue,
        minimum width=2.1cm,
        minimum height=1.55cm,
        align=center
    }
]

% Flip-Flops
\node[ff] (ff3) at (0,0) {
    \textbf{\textenglish{D FF}}\\[1mm]
    \(Q_3\)
};

\node[ff] (ff2) at (3.4,0) {
    \textbf{\textenglish{D FF}}\\[1mm]
    \(Q_2\)
};

\node[ff] (ff1) at (6.8,0) {
    \textbf{\textenglish{D FF}}\\[1mm]
    \(Q_1\)
};

\node[ff] (ff0) at (10.2,0) {
    \textbf{\textenglish{D FF}}\\[1mm]
    \(Q_0\)
};

% Serial input
\draw[->] (-2.0,0) -- (ff3.west);
\node[above] at (-1.15,0.1)
{\textenglish{Serial Input}};

% Data links
\draw[->] (ff3.east) -- (ff2.west);
\draw[->] (ff2.east) -- (ff1.west);
\draw[->] (ff1.east) -- (ff0.west);

% Serial output
\draw[->] (ff0.east) -- ++(1.8,0);
\node[above] at (12.0,0.1)
{\textenglish{Serial Output}};

% Clock line
\draw (-1.4,-2.0) -- (11.0,-2.0);
\node[left] at (-1.4,-2.0) {\textenglish{CLK}};

\foreach \x in {0,3.4,6.8,10.2}
{
    \draw (\x,-2.0) -- (\x,-0.78);
    \fill (\x,-2.0) circle (1.7pt);
}

\end{tikzpicture}
\end{chThreeDiagram}
\end{center}

\par\begingroup

প্রতিটি \textenglish{Clock Pulse}-এ একটি
\textenglish{Flip-Flop}-এর বিট পরবর্তী
\textenglish{Flip-Flop}-এ স্থানান্তরিত হয়।
\par\endgroup



% =========================================================
% Right Shift Example
% =========================================================

\subsubsection*{রাইট শিফটের উদাহরণ
(\textenglish{Right Shift Example})}

ধরা যাক, একটি \(4\)-\textenglish{bit Register}-এ—

\[
Q_3Q_2Q_1Q_0=1011
\]

সংরক্ষিত আছে।

এখন বাম প্রান্ত দিয়ে নতুন বিট \(0\) প্রবেশ করছে এবং
একটি \textenglish{Clock Pulse}-এ তথ্য এক ধাপ ডান দিকে
\textenglish{Shift} হচ্ছে।

\[
1011
\longrightarrow
0101
\]

নিচের চিত্রে বিষয়টি দেখানো হলো—

\begin{center}
\begin{chThreeDiagram}
\begin{tikzpicture}[
    >=Stealth,
    line width=0.8pt,
    bit/.style={
        draw,
        rounded corners=2pt,
        fill=ICTLightGreen,
        minimum width=1.25cm,
        minimum height=0.9cm,
        align=center
    }
]

% Before
\node at (-2.0,1.2) {\textbf{আগে}};

\node[bit] (b3) at (0,1.2) {\Large \(1\)};
\node[bit] (b2) at (1.45,1.2) {\Large \(0\)};
\node[bit] (b1) at (2.90,1.2) {\Large \(1\)};
\node[bit] (b0) at (4.35,1.2) {\Large \(1\)};

\node[above,font=\small] at (0,1.78) {\(Q_3\)};
\node[above,font=\small] at (1.45,1.78) {\(Q_2\)};
\node[above,font=\small] at (2.90,1.78) {\(Q_1\)};
\node[above,font=\small] at (4.35,1.78) {\(Q_0\)};

% New bit
\node at (-1.15,1.2) {\(0\)};
\draw[->] (-0.85,1.2) -- (-0.65,1.2);

\node[above,font=\scriptsize] at (-1.15,1.62)
{\textenglish{New Bit}};

% Clock arrow
\draw[->] (5.15,1.2) -- (6.85,1.2);
\node[above,font=\small] at (6.0,1.35)
{\textenglish{1 Clock Pulse}};

% After
\node at (7.6,1.2) {\textbf{পরে}};

\node[bit] at (8.85,1.2) {\Large \(0\)};
\node[bit] at (10.30,1.2) {\Large \(1\)};
\node[bit] at (11.75,1.2) {\Large \(0\)};
\node[bit] at (13.20,1.2) {\Large \(1\)};

\node[above,font=\small] at (8.85,1.78) {\(Q_3\)};
\node[above,font=\small] at (10.30,1.78) {\(Q_2\)};
\node[above,font=\small] at (11.75,1.78) {\(Q_1\)};
\node[above,font=\small] at (13.20,1.78) {\(Q_0\)};

\end{tikzpicture}
\end{chThreeDiagram}
\end{center}

এখানে—

\[
Q_3\leftarrow0
\]

\[
Q_2\leftarrow Q_3
\]

\[
Q_1\leftarrow Q_2
\]

\[
Q_0\leftarrow Q_1
\]

অতএব,

\[
\boxed{
1011
\xrightarrow{\text{\textenglish{1 Clock Pulse}}}
0101
}
\]

\par\begingroup

\textenglish{Shift}-এর দিক এবং নতুন বিট কোন প্রান্ত দিয়ে প্রবেশ করবে,
তা সার্কিটের সংযোগের ওপর নির্ভর করে।
এখানে \textenglish{Right Shift} দেখানো হয়েছে।
\par\endgroup



% =========================================================
% Types of Shift Register
% =========================================================

\subsection{শিফট রেজিস্টারের প্রকারভেদ}

ডেটা কীভাবে প্রবেশ করে এবং কীভাবে বের হয় তার ভিত্তিতে
\textenglish{Shift Register} প্রধানত চার প্রকার—

\begin{center}
\renewcommand{\arraystretch}{1.25}

\begin{tabularx}
{\textwidth}
{|L{0.14\textwidth}|L{0.29\textwidth}|Y|}

\hline

\rowcolor{ICTLightBlue}

\textbf{প্রকার}
&
\textbf{পূর্ণরূপ}
&
\textbf{কাজ}
\\
\hline

\textenglish{SISO}
&
\textenglish{Serial-In Serial-Out}
&
সিরিয়ালভাবে তথ্য গ্রহণ এবং সিরিয়ালভাবে তথ্য প্রদান করে।
\\
\hline

\textenglish{SIPO}
&
\textenglish{Serial-In Parallel-Out}
&
সিরিয়ালভাবে তথ্য গ্রহণ করে প্যারালাল আউটপুট প্রদান করে।
\\
\hline

\textenglish{PISO}
&
\textenglish{Parallel-In Serial-Out}
&
একসঙ্গে একাধিক বিট গ্রহণ করে সিরিয়ালভাবে আউটপুট দেয়।
\\
\hline

\textenglish{PIPO}
&
\textenglish{Parallel-In Parallel-Out}
&
একসঙ্গে একাধিক বিট গ্রহণ ও প্রদান করে।
\\
\hline

\end{tabularx}
\end{center}


% =========================================================
% Shift Register Applications
% =========================================================

\subsubsection*{শিফট রেজিস্টারের ব্যবহার}

\begin{enumerate}[leftmargin=*,itemsep=2pt]

    \item \textenglish{Serial} ও \textenglish{Parallel Data}-এর মধ্যে
    রূপান্তর করতে।

    \item ডিজিটাল যোগাযোগ ব্যবস্থায় তথ্য স্থানান্তর করতে।

    \item বাইনারি তথ্য সাময়িকভাবে সংরক্ষণ করতে।

    \item নির্দিষ্ট ডিজিটাল সার্কিটে তথ্য বিলম্বিত করতে।

    \item কিছু \textenglish{Counter} ও
    \textenglish{Sequence Generator} তৈরিতে।

\end{enumerate}


% =========================================================
% 3.12.6 Buffer Register
% =========================================================

\subsection{বাফার রেজিস্টার
(\textenglish{Buffer Register})}

যে রেজিস্টার ডেটা এক ডিজিটাল অংশ থেকে অন্য অংশে স্থানান্তরের সময়
সাময়িকভাবে ধরে রাখে, তাকে
\textbf{\textenglish{Buffer Register}} বলে।

এটি মূলত তথ্যের একটি অস্থায়ী সংরক্ষণস্থান হিসেবে কাজ করে।

\subsubsection*{কার্যপদ্ধতি}

ডেটা রেজিস্টারে লোড হওয়ার পর পরবর্তী নির্দেশ বা নিয়ন্ত্রণ সংকেত
না আসা পর্যন্ত সেটি সেখানে ধরে রাখা যায়।

এর ফলে দুটি ডিজিটাল ইউনিটের মধ্যে তথ্য আদান-প্রদানের সময়
ডেটা সাময়িকভাবে সংরক্ষণ করা সম্ভব হয়।

\subsubsection*{বাফার রেজিস্টারের ব্যবহার}

\begin{enumerate}[leftmargin=*,itemsep=2pt]

    \item \textenglish{CPU} ও অন্যান্য ডিজিটাল ইউনিটের মধ্যে
    ডেটা সাময়িকভাবে ধরে রাখতে।

    \item ডেটা স্থানান্তরের সময় দুই সার্কিটের কাজের সময়ের পার্থক্য
    সামলাতে।

    \item ইনপুট বা আউটপুট ডেটা সাময়িকভাবে সংরক্ষণ করতে।

    \item প্রসেসিংয়ের আগে বা পরে ডেটা ধরে রাখতে।

\end{enumerate}


% =========================================================
% 3.12.7 Parallel Load Register
% =========================================================

\subsection{প্যারালাল লোড রেজিস্টার
(\textenglish{Parallel Load Register})}

যে রেজিস্টারে একাধিক বিটের বাইনারি তথ্য একই সময়ে পৃথক ইনপুটের মাধ্যমে
গ্রহণ ও সংরক্ষণ করা যায়, তাকে
\textbf{\textenglish{Parallel Load Register}} বলে।

একটি \(4\)-\textenglish{bit Parallel Load Register}-এ
চারটি \textenglish{Flip-Flop} থাকে এবং প্রতিটি বিটের জন্য
একটি করে আলাদা ডেটা ইনপুট থাকে।

একই সক্রিয় \textenglish{Clock}-এ চারটি ইনপুট বিট
একসঙ্গে সংশ্লিষ্ট \textenglish{Flip-Flop}-এ সংরক্ষিত হয়।


% =========================================================
% Parallel Load Diagram
% =========================================================

\subsubsection*{৪-বিট প্যারালাল লোড রেজিস্টারের গঠন}

\begin{center}
\begin{chThreeDiagram}
\begin{tikzpicture}[
    >=Stealth,
    line width=0.8pt,
    ff/.style={
        draw,
        rounded corners=2pt,
        fill=ICTLightGreen,
        minimum width=2.05cm,
        minimum height=1.55cm,
        align=center
    }
]

% Flip-Flops
\node[ff] (ff3) at (0,0) {
    \textbf{\textenglish{D FF}}\\[1mm]
    \(Q_3\)
};

\node[ff] (ff2) at (3.4,0) {
    \textbf{\textenglish{D FF}}\\[1mm]
    \(Q_2\)
};

\node[ff] (ff1) at (6.8,0) {
    \textbf{\textenglish{D FF}}\\[1mm]
    \(Q_1\)
};

\node[ff] (ff0) at (10.2,0) {
    \textbf{\textenglish{D FF}}\\[1mm]
    \(Q_0\)
};

% Data inputs
\draw[->] (0,2.0) -- (0,0.78);
\node[above] at (0,2.0) {\(D_3\)};

\draw[->] (3.4,2.0) -- (3.4,0.78);
\node[above] at (3.4,2.0) {\(D_2\)};

\draw[->] (6.8,2.0) -- (6.8,0.78);
\node[above] at (6.8,2.0) {\(D_1\)};

\draw[->] (10.2,2.0) -- (10.2,0.78);
\node[above] at (10.2,2.0) {\(D_0\)};

% Outputs
\draw[->] (ff3.south) -- ++(0,-1.1);
\node[below] at (0,-1.1) {\(Q_3\)};

\draw[->] (ff2.south) -- ++(0,-1.1);
\node[below] at (3.4,-1.1) {\(Q_2\)};

\draw[->] (ff1.south) -- ++(0,-1.1);
\node[below] at (6.8,-1.1) {\(Q_1\)};

\draw[->] (ff0.south) -- ++(0,-1.1);
\node[below] at (10.2,-1.1) {\(Q_0\)};

% Common clock
\draw (-1.2,-2.0) -- (11.2,-2.0);
\node[left] at (-1.2,-2.0) {\textenglish{CLK}};

\foreach \x in {0,3.4,6.8,10.2}
{
    \draw (\x,-2.0) -- (\x,-0.78);
    \fill (\x,-2.0) circle (1.7pt);
}

\end{tikzpicture}
\end{chThreeDiagram}
\end{center}

এখানে,

\[
D_3D_2D_1D_0
\]

হলো চারটি প্যারালাল ইনপুট এবং,

\[
Q_3Q_2Q_1Q_0
\]

হলো চারটি সংরক্ষিত বিট।

যদি,

\[
D_3D_2D_1D_0=1101
\]

হয়, তাহলে সক্রিয় \textenglish{Clock}-এর পর—

\[
Q_3Q_2Q_1Q_0=1101
\]

হবে।

অর্থাৎ,

\[
D_3=1\Rightarrow Q_3=1
\]

\[
D_2=1\Rightarrow Q_2=1
\]

\[
D_1=0\Rightarrow Q_1=0
\]

\[
D_0=1\Rightarrow Q_0=1
\]

সুতরাং,

\[
\boxed{
1101
\xrightarrow{\text{\textenglish{Parallel Load}}}
1101
}
\]


% =========================================================
% Buffer vs Parallel Load
% =========================================================

\subsection{বাফার ও প্যারালাল লোড রেজিস্টারের সম্পর্ক}

\begin{center}
\renewcommand{\arraystretch}{1.25}

\begin{tabularx}
{\textwidth}
{|L{0.21\textwidth}|Y|Y|}

\hline

\rowcolor{ICTLightBlue}

\textbf{বিষয়}
&
\textbf{\textenglish{Buffer Register}}
&
\textbf{\textenglish{Parallel Load Register}}
\\
\hline

প্রধান উদ্দেশ্য
&
ডেটা সাময়িকভাবে ধরে রাখা।
&
একাধিক বিট একই সময়ে সংরক্ষণ করা।
\\
\hline

ডেটা গ্রহণ
&
সার্কিটের নকশা অনুযায়ী।
&
একাধিক ইনপুট থেকে একই সময়ে।
\\
\hline

নিয়ন্ত্রণ
&
\textenglish{Load} বা \textenglish{Enable} সংকেত থাকতে পারে।
&
প্যারালাল ডেটা গ্রহণের ব্যবস্থা থাকে।
\\
\hline

ব্যবহার
&
ডেটা স্থানান্তরের মধ্যবর্তী ধাপে।
&
দ্রুত একাধিক বিট সংরক্ষণে।
\\
\hline

\end{tabularx}
\end{center}

\par\begingroup

\textbf{মনে রাখবে:}
\textenglish{Buffer Register} এবং
\textenglish{Parallel Load Register} সম্পূর্ণ আলাদা শ্রেণি নয়।
একটি \textenglish{Buffer Register}-ও প্যারালালভাবে ডেটা লোড করার
সুবিধাসহ তৈরি হতে পারে।
\par\endgroup

% Each drawing represents a complete data path; clock branches are independent.
\newcommand{\chThreeSerialRegister}[1]{%
\begin{center}
\begin{minipage}{\linewidth}\centering\begin{chThreeDiagram}\begin{tikzpicture}[>=Stealth,line width=.7pt]
\foreach \i/\b in {0/3,1/2,2/1,3/0}{
  \node[chff] (f\i) at (2.7*\i,0) {D FF\\$Q_\b$};
  \draw (f\i.south)--++(0,-.8);
  \node[font=\scriptsize] at (2.7*\i,-.35) {$\triangleright$};
  \fill (2.7*\i,-1.6) circle(1.4pt);
  \ifnum#1=1\draw (f\i.east)--++(.25,0)--++(0,1.3) node[above] {$Q_\b$};\fi
}
\draw (-2,0)--(f0.west) node[pos=0,left] {SI};
\foreach \i/\j in {0/1,1/2,2/3}{\draw[->] (f\i.east)--(f\j.west);}
\draw[->] (f3.east)--++(1.5,0) node[right] {SO};
\draw (-1,-1.6)--(8.1,-1.6);\node[left] at (-1,-1.6) {CLK};
\end{tikzpicture}\end{chThreeDiagram}
\ifnum#1=1\captionof{figure}{প্রতিটি Q থেকে আলাদা আউটপুটসহ চার বিটের SIPO।}
\else\captionof{figure}{চার D flip-flop-এর ধারাবাহিক সংযোগে SISO।}\fi

\end{minipage}
\end{center}}

\subsubsection*{SISO: সিরিয়াল ইনপুট, সিরিয়াল আউটপুট}
\chThreeSerialRegister{0}
\noindent\textbf{সার্কিটের ব্যাখ্যা:} চার D flip-flop-এর একটি Q আউটপুট পরের D ইনপুটে যুক্ত। সবগুলো একই সক্রিয় ক্লক পায়। SI থেকে প্রতি ক্লকে একটি বিট ঢোকে; পুরোনো বিটগুলো এক ধাপ ডানে সরে। SO হলো শেষ flip-flop-এর Q। একই ক্লকে সবাই \emph{আগের} Q গ্রহণ করে; একটি বিট একই ক্লকে সব ধাপ পার হয় না।

\chThreeExample{১}{SISO-তে একটি বিটের যাত্রা}
শুরু $Q_3Q_2Q_1Q_0=0000$; SI-তে প্রথমে ১, তারপর তিনবার ০।
\chThreeAnswer $0000\to1000\to0100\to0010\to0001$। চার নম্বর সক্রিয় ক্লকের পর SO=১। SO-কে সর্বশেষ Q হিসেবে পড়া হয়েছে; বেরিয়ে যাওয়া \emph{পুরোনো} বিট পড়ার নিয়ম আলাদা করে উল্লেখ করতে হয়।

\subsubsection*{SIPO: সিরিয়াল ইনপুট, প্যারালাল আউটপুট}
\chThreeSerialRegister{1}
\noindent\textbf{সার্কিটের ব্যাখ্যা:} ইনপুট পথ SISO-এর মতো, কিন্তু প্রতিটি Q থেকে আলাদা output tap নেওয়া হয়েছে। এজন্য সিরিয়ালভাবে ঢোকা চার বিট শেষে একসঙ্গে পড়া যায়। LED-তে একাধিক বিট পাঠানো বা serial data-কে parallel data-য় বদলাতে এটি কাজে লাগে।

\chThreeExample{২}{SIPO-তে ১০১১ লোড}
SI থেকে ক্রমান্বয়ে ১, ১, ০, ১ পাঠাও; শুরুতে 0000।
\chThreeAnswer $0000\to1000\to1100\to0110\to1011$। এই ডান শিফটের সংযোগে কাঙ্ক্ষিত 1011 পেতে সর্বডানের বিট আগে পাঠানো হয়েছে। কোন প্রান্তে ইনপুট ঢোকে তা না জানলে bit order ঠিক করা যায় না।

\subsubsection*{PISO: প্যারালাল ইনপুট, সিরিয়াল আউটপুট}
\begin{center}
\begin{minipage}{\linewidth}\centering\begin{chThreeDiagram}\begin{tikzpicture}[>=Stealth,line width=.7pt]
\foreach \i/\b in {0/3,1/2,2/1,3/0}{
 \node[chff] (f\i) at (3.2*\i,0) {D FF\\$Q_\b$};
 \node[chmux] (m\i) at (3.2*\i,2) {MUX\\$L$};
 \draw[->] (m\i.south)--(f\i.north) node[midway,right,font=\scriptsize] {D};
 \draw[->] (3.2*\i,3.3)--(m\i.north) node[pos=0,above] {$P_\b$};
 \draw (f\i.south)--(3.2*\i,-1.6);
 \fill (3.2*\i,-1.6) circle(1.4pt);
}
\draw[->] (-1.6,2)--(m0.west) node[pos=0,left] {SI};
\foreach \i/\j in {0/1,1/2,2/3}{\draw[->] (f\i.east)--++(.55,0)|-(m\j.west);}
\draw[->] (f3.east)--++(1,0) node[right] {SO};
\draw (-1,-1.6)--(9.6,-1.6) node[pos=0,left] {CLK};
\node[font=\small] at (4.8,4.1) {সব MUX-এ একই $L$: ১=লোড, ০=শিফট};
\end{tikzpicture}\end{chThreeDiagram}
\captionof{figure}{PISO-তে MUX প্যারালাল ডেটা অথবা আগের ধাপের Q নির্বাচন করে।}

\end{minipage}
\end{center}
\noindent\textbf{সার্কিটের ব্যাখ্যা:} প্রতিটি D-এর আগে একটি দুই-পথের selector বা multiplexer আছে। $L=1$ হলে চারটি P ইনপুট পরের সক্রিয় ক্লকে একসঙ্গে লোড হয়। $L=0$ হলে প্রতিটি ধাপ আগের ধাপের Q গ্রহণ করে ডানে শিফট হয়। $D_3=LP_3+\overline L\,SI$ এবং $D_i=LP_i+\overline LQ_{i+1}$, $i=0,1,2$।

\chThreeExample{৩}{PISO থেকে ১১০১ পাঠানো}
লোড ক্লকের পরে Q=1101। SO থেকে বর্তমান $Q_0$ পড়ে তারপর SI=০ দিয়ে শিফট করো।
\chThreeAnswer পাঠানো বিটের ক্রম ১, ০, ১, ১। শিফটের পর state যথাক্রমে 0110, 0011, 0001, 0000। এখানে পড়া হয়েছে শিফটের \emph{আগে}; তাই SISO উদাহরণের ``ক্লকের পরে SO'' নিয়মের সঙ্গে গুলিয়ে ফেলা যাবে না।

\subsubsection*{PIPO: প্যারালাল ইনপুট, প্যারালাল আউটপুট}
\begin{center}
\begin{minipage}{\linewidth}\centering\begin{chThreeDiagram}\begin{tikzpicture}[>=Stealth,line width=.7pt]
\foreach \i/\b in {0/3,1/2,2/1,3/0}{
 \node[chff] (f\i) at (2.8*\i,0) {D FF\\$Q_\b$};
 \draw[->] (2.8*\i,1.7)--(f\i.north) node[pos=0,above] {$P_\b$};
 \draw[->] (f\i.east)--++(.6,0) node[right] {$Q_\b$};
 \draw (f\i.south)--(2.8*\i,-1.6);\fill (2.8*\i,-1.6) circle(1.4pt);
}
\draw (-1,-1.6)--(8.4,-1.6) node[pos=0,left] {CLK};
\end{tikzpicture}\end{chThreeDiagram}
\captionof{figure}{চার পৃথক D ইনপুট একই ক্লকে সংরক্ষিত হয়ে চার Q আউটপুট দেয়।}

\end{minipage}
\end{center}
\noindent\textbf{সার্কিটের ব্যাখ্যা:} প্রতিটি ইনপুট তার নিজস্ব D flip-flop-এ যায়; এক Q থেকে পরের D-তে শিফট সংযোগ নেই। এক সক্রিয় ক্লকে চার বিট লোড হয়, এবং চার আউটপুট একসঙ্গে পাওয়া যায়। এটি প্যারালাল লোড রেজিস্টারের সরল নির্মাণ।

\chThreeExample{৪}{কত flip-flop, কত ক্লক?}
আট বিট সংরক্ষণে আটটি flip-flop লাগে। SIPO-তে আট বিট ভর্তি করতে আট সক্রিয় ক্লক লাগে; PIPO-তে একটি লোড ক্লক যথেষ্ট।
\chThreeAnswer দুটি রেজিস্টারই আট বিট রাখে; ডেটা গ্রহণের পথ আলাদা। আট flip-flop মানে আটটি সংখ্যা নয়, বরং $2^8=256$টি সম্ভাব্য আট-বিট pattern-এর যেকোনো একটি।

% Mahbub teaching-type follow-up, 2026-10-14
\chThreeExample{৫}{প্রাথমিক ডেটা বের হওয়া ও নতুন ডেটা ঢোকা}
উপরের SISO-তে $Q_3Q_2Q_1Q_0=1010$। পরপর SI-তে ০, ১, ১, ০ দাও। প্রতিটি clock-এর আগে SO পড়লে পুরোনো ডেটা কোন ক্রমে বের হবে এবং শেষ state কী?
\chThreeAnswer SO হলো $Q_0$। চার pulse-এ state $1010\to0101\to1010\to1101\to0110$। clock-এর আগে পড়া SO-এর ক্রম ০, ১, ০, ১: পুরোনো শব্দ LSB আগে বের হয়েছে। নতুন input ০, ১, ১, ০ শেষে state-এ বিপরীত অবস্থানক্রমে বসেছে। পুরোনো output-এর sampling এবং clock-এর পরের state একই ধাপের দুই আলাদা পর্যবেক্ষণ।


\subsection{বাফার, দুই দিকে শিফট এবং universal register}
\smallhead{বাফার রেজিস্টারে load ও hold}
শুধু ডেটা লোড করলেই হয় না; পরবর্তী অংশ প্রস্তুত না হওয়া পর্যন্ত সেটি ধরে রাখতে হয়। Load=১-এ নতুন P গ্রহণ এবং Load=০-এ নিজের Q ধরে রাখতে প্রতিটি D-এর আগে selector বসানো যায়।
\begin{center}
\begin{minipage}{\linewidth}\centering\begin{chThreeDiagram}\begin{tikzpicture}[>=Stealth,line width=.7pt]
\node[chmux] (m) at (0,0) {MUX\\$L$};
\node[chff] (f) at (3,0) {D FF\\$Q_i$};
\draw[->] (-1.6,.35)--(-.8,.35)--(m.west) node[pos=0,left] {$P_i$};
\draw[->] (m.east)--(f.west) node[midway,above] {$D_i$};
\draw[->] (f.east)--(5,0) node[right] {$Q_i$};
\draw[->] (4.3,0)--(4.3,1.5)--(-1,1.5)|-(m.north);
\fill (4.3,0) circle(1.5pt);
\draw[->] (3,-1.7)--(f.south) node[pos=0,below] {CLK};
\end{tikzpicture}\end{chThreeDiagram}
\captionof{figure}{বাফার রেজিস্টারের একটি বিটে নতুন ডেটা ও নিজের পুরোনো Q-এর নির্বাচন।}

\end{minipage}
\end{center}
সমীকরণ $D_i=LP_i+\overline LQ_i$। আট বিটে এই অংশ আটবার থাকে এবং ক্লক ও L সাধারণভাবে সরবরাহ করা হয়। বাফার বলা হয়েছে কাজের উদ্দেশ্য অনুযায়ী; এটি PIPO নির্মাণও হতে পারে।

\chThreeExample{৬}{লোড না থাকলে ডেটা ধরে রাখা}
বর্তমান Q=1010, ইনপুট P=0111।
\chThreeAnswer $L=0$-এ ক্লক এলেও Q=1010; $L=1$-এ পরের সক্রিয় ক্লকে Q=0111। ইনপুট বদলানো এবং stored output বদলানো এক ঘটনা নয়।

\smallhead{দুই দিকে শিফট রেজিস্টারের একটি বিট}
\begin{center}
\begin{minipage}{\linewidth}\centering\begin{chThreeDiagram}\begin{tikzpicture}[>=Stealth,line width=.7pt]
\node[chmux] (m) at (0,0) {MUX\\$R$};
\node[chff] (f) at (3,0) {D FF\\$Q_i$};
\draw[->] (-2,.5)--(-.8,.5)--(m.west) node[pos=0,left] {$Q_{i+1}$};
\draw[->] (-2,-.5)--(-.8,-.5)--(m.south) node[pos=0,left] {$Q_{i-1}$};
\draw[->] (m.east)--(f.west) node[midway,above] {$D_i$};
\draw[->] (f.east)--++(1,0) node[right] {$Q_i$};
\draw[->] (3,-1.7)--(f.south) node[pos=0,below] {CLK};
\end{tikzpicture}\end{chThreeDiagram}
\captionof{figure}{দিক নিয়ন্ত্রণ R অনুযায়ী পাশের দুই flip-flop-এর একটি নির্বাচন।}

\end{minipage}
\end{center}
$R=1$-এ ডান শিফট: $D_i=Q_{i+1}$। $R=0$-এ বাম শিফট: $D_i=Q_{i-1}$। দুই প্রান্তে প্রতিবেশীর বদলে যথাক্রমে ডান-শিফট ও বাম-শিফট serial input দেওয়া হয়। তাই সাধারণ মধ্যবর্তী বিটে $D_i=RQ_{i+1}+\overline RQ_{i-1}$।

\chThreeExample{৭}{বাম ও ডান শিফটের মান}
চার বিট 0110-তে প্রান্তিক নতুন বিট ০।
\chThreeAnswer ডান শিফটে 0011, অর্থাৎ $6\to3$; বাম শিফটে 1100, অর্থাৎ $6\to12$। চিহ্নবিহীন সংখ্যায় logical right shift দুই দিয়ে পূর্ণসংখ্যা ভাগের মতো। বাম শিফটে বিট বেরিয়ে গেলে overflow হতে পারে; যেমন $1100\to1000$ প্রকৃত $12\times2$ নয়। signed arithmetic right shift-এ sign bit ধরে রাখার নিয়ম আলাদা।

\smallhead{Universal shift register}
একটি চার-পথের selector দিয়ে প্রতিটি D-তে নিজের Q, ডান পাশের/বাম পাশের প্রতিবেশী এবং parallel input-এর যেকোনো একটি নেওয়া যায়। এতে hold, ডান শিফট, বাম শিফট ও parallel load—চার কাজ হয়।
\begin{center}
\begin{minipage}{\linewidth}\centering\begin{chThreeDiagram}\begin{tikzpicture}[>=Stealth,line width=.7pt]
\node[draw,minimum width=23mm,minimum height=30mm,align=center] (m) at (0,0) {4:1 MUX\\$S_1S_0$};
\node[chff] (f) at (4,0) {D FF\\$Q_i$};
\foreach \y/\label in {1.1/Q_i,.4/Q_{i+1},-.4/Q_{i-1},-1.1/P_i}{\draw[->] (-3,\y)--(-1.15,\y) node[pos=0,left] {$\label$};}
\draw[->] (m.east)--(f.west) node[midway,above] {$D_i$};
\draw[->] (f.east)--++(1,0) node[right] {$Q_i$};
\draw[->] (4,-2.2)--(f.south) node[pos=0,below] {CLK};
\end{tikzpicture}\end{chThreeDiagram}
\captionof{figure}{Universal register-এর প্রতিটি বিটে চারটি ডেটা পথের নির্বাচন।}

\end{minipage}
\end{center}
এই উদাহরণের control mapping: 00=hold, 01=ডান শিফট, 10=বাম শিফট, 11=parallel load। পুরো রেজিস্টারে একাধিক এমন অংশ, সাধারণ ক্লক ও সাধারণ selector bits থাকে। বাস্তব IC-এর mapping তার datasheet অনুযায়ী নিতে হয়।

